\begin{table}[H]
\centering
\caption{Comparação entre especificações GMM do sistema AIDS}
\label{tab:comparacao-modelos}
\small % <-- Tamanho intermediário
\setlength{\tabcolsep}{6pt} % <-- Espaçamento padrão
\renewcommand{\arraystretch}{1.15} % <-- Altura das linhas moderada
\begin{tabular}{@{}p{7.5cm}rrrrr@{}} % <-- Largura da 1ª coluna ajustada
\toprule
Especificação 
& \(N\) 
& Parâmetros 
& Momentos 
& gl overid. 
& \(J\) \\
\midrule
Modelo sem restrições teóricas
& 40 & 18 & 18 & 0 & 0.0000 \\

Modelo com homogeneidade
& 40 & 15 & 18 & 3 & 14.8150 \\

Modelo com homogeneidade e simetria
& 40 & 12 & 18 & 6 & 36.4051 \\

Modelo com homogeneidade, simetria e instrumentos ampliados
& 39 & 12 & 30 & 18 & 60.8732 \\
\bottomrule
\end{tabular}

\par\vspace{0.35em}
\begin{minipage}{0.95\textwidth} % <-- Minipage com 95% da largura
\scriptsize % <-- Nota ligeiramente menor que a tabela
\textit{Nota:} A coluna \(J\) reporta a estatística de sobreidentificação do GMM. O primeiro modelo é exatamente identificado, pois o número de momentos é igual ao número de parâmetros. Nos demais modelos, as restrições teóricas e/ou a ampliação do conjunto de instrumentos geram graus de liberdade de sobreidentificação.
\end{minipage}
\end{table}